\section{Filtros morfológicos}

    \subsection{Subtitulo}

        Desarrollar el ejercicio por subsecciones queda prolijo \dots

    \subsection{Preguntas y análisis}

        \begin{itemize}[leftmargin=*]
            \item \textbf{¿Qué efecto tuvo la apertura morfológica sobre la máscara? ¿Se eliminó el ruido externo? ¿Se afectaron los bordes de los objetos de interés? }
            
                Desarrollar la respuesta acá \dots

            \item \textbf{¿Qué efecto tuvo el cierre morfológico? ¿Se rellenaron los huecos internos? ¿Se unieron objetos que estaban separados?} 
            
                Desarrollar la respuesta acá \dots

            \item \textbf{¿Cómo afecta el tamaño del elemento estructurante a los resultados? ¿Existe un tamaño óptimo según el tamaño de los objetos y del ruido? }
            
                Desarrollar la respuesta acá \dots

            \item \textbf{¿La forma del elemento estructurante (cuadrado vs. disco vs. cruz) produce diferencias notables? ¿En qué contextos podría ser relevante?} 
            
                Desarrollar la respuesta acá \dots

            \item \textbf{Probar la robustez del pipeline completo (segmentación + morfología) con otras imágenes de la misma clase, pero capturadas bajo condiciones diferentes (iluminación distinta, cámara diferente, fondo variable). ¿Los mismos parámetros siguen funcionando bien? ¿Qué ajustes serían necesarios?} 
            
                Desarrollar la respuesta acá \dots

            \item \textbf{¿Es suficiente aplicar solo apertura o solo cierre, o es necesario encadenar ambas operaciones? ¿En qué orden? Justificar según el tipo de ruido presente.} 
            
                Desarrollar la respuesta acá \dots

            \item \textbf{En un escenario de inspección industrial automatizada, ¿cuán crítica es la selección de estos parámetros morfológicos para garantizar una detección confiable de defectos o conteo de objetos?} 
            
                Desarrollar la respuesta acá \dots
                
        \end{itemize}
